\documentclass[11pt]{article}

\usepackage[T1]{fontenc}
\usepackage[utf8]{inputenc}
\usepackage{lmodern}
\usepackage{microtype}
\usepackage{amsmath,amssymb,amsthm,mathtools}
\usepackage{aliascnt}
\usepackage{enumitem}
\usepackage{xcolor}
\usepackage{booktabs}
\usepackage{array}
\usepackage{longtable}
\usepackage{xurl}
\usepackage{hyperref}
\usepackage[nameinlink,capitalise,noabbrev]{cleveref}

\emergencystretch=6em

\hypersetup{
  colorlinks=true,
  linkcolor=blue!45!black,
  citecolor=blue!45!black,
  urlcolor=blue!45!black
}

\newtheorem{theorem}{Theorem}[section]
\newaliascnt{lemma}{theorem}
\newtheorem{lemma}[lemma]{Lemma}
\aliascntresetthe{lemma}
\newaliascnt{proposition}{theorem}
\newtheorem{proposition}[proposition]{Proposition}
\aliascntresetthe{proposition}
\newaliascnt{corollary}{theorem}
\newtheorem{corollary}[corollary]{Corollary}
\aliascntresetthe{corollary}
\newaliascnt{definition}{theorem}
\newtheorem{definition}[definition]{Definition}
\aliascntresetthe{definition}
\newaliascnt{conjecture}{theorem}
\newtheorem{conjecture}[conjecture]{Conjecture}
\aliascntresetthe{conjecture}
\newaliascnt{problem}{theorem}
\newtheorem{problem}[problem]{Problem}
\aliascntresetthe{problem}
\newaliascnt{remark}{theorem}
\newtheorem{remark}[remark]{Remark}
\aliascntresetthe{remark}
\crefname{lemma}{Lemma}{Lemmas}
\crefname{proposition}{Proposition}{Propositions}
\crefname{corollary}{Corollary}{Corollaries}
\crefname{definition}{Definition}{Definitions}
\crefname{conjecture}{Conjecture}{Conjectures}
\crefname{problem}{Problem}{Problems}
\crefname{remark}{Remark}{Remarks}

\newcommand{\F}{\mathbb{F}}
\newcommand{\RS}{\operatorname{RS}}
\newcommand{\MCA}{\operatorname{MCA}}
\newcommand{\dist}{\operatorname{dist}}
\newcommand{\image}{\operatorname{im}}
\newcommand{\Bad}{\operatorname{Bad}}
\newcommand{\Ball}{\operatorname{Ball}}
\newcommand{\List}{\Lambda}
\newcommand{\epsstar}{\varepsilon^\ast}
\newcommand{\deltastar}{\delta^\ast}
\newcommand{\deltacap}{\delta_{\mathrm{cap}}}
\newcommand{\Cpair}{\mathcal C^{\equiv 2}}
\newcommand{\Lean}{\textsc{Lean}}
\newcommand{\lean}[1]{\begingroup\Urlmuskip=0mu plus 1mu\path{#1}\endgroup}
\newenvironment{leanledger}{\small\raggedright}{\par}
\newcommand{\todo}[1]{\textcolor{red!65!black}{[TODO: #1]}}


\title{Capacity-Radius Proximity for Smooth-Domain Reed--Solomon Codes}
\author{Chloe van der Vlugt}
\date{}

\hypersetup{
  pdftitle={Capacity-Radius Proximity for Smooth-Domain Reed--Solomon Codes},
  pdfauthor={Chloe van der Vlugt}
}

\begin{document}

\maketitle

\begin{abstract}
We study list decoding and same-set mutual correlated agreement for
Reed--Solomon codes over smooth evaluation domains. A two-coordinate plane
is exactly an interleaved pair word, equivalently the base--direction pair
of a scalar affine line.

At count radius \(n-k\), an agreement-set injection and a matching
received-word construction prove
\[
  L_m\!\left(\mathcal C,\frac{n-k}{n}\right)=\binom nk
\]
for \(1\le k\le n\) and every positive interleaving width. A next-shell
construction gives the three-regime fractional-list phase split. For pair
mutual correlated agreement, compressed interpolation certificates give
\[
  |\Bad_{\MCA}|\le\binom n{k+1}
\]
at every radius. A smooth rate-\(1/4\) family over \(\mathbb F_{3^n}\)
meets both \(2^{-128}\) fractional budgets: its list target holds at capacity
and its pair-MCA target holds at every normalized radius.
For the separate prime-field instance
\((q,n,k)=(2^{128}+192513,4096,1024)\), the exact bad-count spectrum has
\(B_{\mathcal C}(0)=1\) and \(B_{\mathcal C}(1)\ge2\). Its greatest safe MCA
count radius is zero, its first unsafe count is one, and its real safe set
is \([0,1/4096)\). The intermediate field-size list optimum remains open.

For the shared systematic cellular cochain construction, we prove a
quantitative comparison for its existing committed-cochain projection and
intrinsic modular coboundary validity.
Its independently defined native distances satisfy
\(d_C\le d_P=D_2\le2d_C\), and a well-formed raw payload in the canonical
\(q=12289\) instance attains \(3/4\) at all three levels. This application
uses a separate field from the two fractional-budget instances above.
Euclidean and full-HCP source comparisons remain separate.
\end{abstract}

\tableofcontents

\section{Introduction}

The Grand List Decoding and Grand Mutual Correlated Agreement challenges ask
for quantitative proximity guarantees for constant-rate Reed--Solomon codes
over smooth evaluation domains \cite{arnonBonehFenzi2026}.  We work with the
two-fold interleaving
\(\Cpair\subseteq(\F^2)^n\).  Every received pair
\(p=(p_0,p_1)\) determines the scalar affine line
\[
  p\in(\F^2)^n
  \quad\longmapsto\quad
  y_\gamma=p_0+\gamma p_1\in\F^n .
\]
List decoding controls the ambiguity of the full pair, while mutual correlated
agreement bounds the fraction of challenges for which scalar agreement fails
to extend to pair agreement on the same coordinate set.

At the capacity-grid radius \((n-k)/n\), we prove the exact universal list
size
\[
  L_m\!\left(\mathcal C,\frac{n-k}{n}\right)=\binom nk
\]
for every positive interleaving width.  A matching next-shell construction
gives the resulting three-regime field-size phase split. For pair MCA, an
injective compressed-certificate map proves the radius-independent bound
\[
  |\Bad_{\MCA}|\le\binom n{k+1}.
\]
An explicit smooth rate-\(1/4\) family over \(\mathbb F_{3^n}\) meets both
\(2^{-128}\) fractional targets. Its actual field/subgroup realization,
endpoint constructions, MCA certificate bound, complete phase split, real
safe intervals and generic distance/calibration endpoints are machine-checked
in Lean. The endpoint arguments are elementary interpolation; the exponential
alphabet and allowable lists do not establish an efficient decoder or a
polynomial-list improvement.

Recent limitations clarify this distinction. Alrabiah, Goyal and Guruswami
\cite[Theorems~1.1, 1.2 and~1.4]{alrabiahGoyalGuruswami2026} construct
super-polynomial lists just beyond Johnson for certain full-length
low-characteristic codes, and growing polynomial lists below capacity for
specified full-length prime-field codes. Their full length is the field
cardinality. Our family instead uses \(n\) subgroup points in
\(\mathbb F_{3^n}\), and its endpoint list \(\binom nk\) is allowed to
grow exponentially in \(n\) while meeting a fraction-of-field budget.
Neither that budget nor the separate pair-MCA event asserts polynomial-list
decoding up to capacity.

To resolve the first nontrivial MCA transition rather than only give a
uniform sufficient bound, we define the exact spectrum
\[
  B_{\mathcal C}(r)=\max_p
  |\Bad_{\MCA}(\mathcal C,r/n,p)|.
\]
For the certified smooth prime-field instance
\(q=2^{128}+192513\equiv1\pmod {4096}\), \(n=4096\), and \(k=1024\), we prove
\(B_{\mathcal C}(0)=1\) and \(B_{\mathcal C}(1)\ge2\).  It follows that zero
is the greatest safe count radius and one the first unsafe count radius at
the \(2^{-128}\) MCA error target. Over real radii the safe interval is
\([0,1/4096)\), so its supremum is \(1/4096\) and is not itself safe.

The coding theorems quantify over the complete ambient received-word and
affine-line spaces.  Their \(2^{-128}\) instantiation uses the displayed
exponential-size field and is not a claim about every supplied smooth code.
An application with an independently defined source metric additionally
requires a comparison with the extracted pair-code distance.

The source application uses the cochain construction and observation profile
of \cite{vanDerVlugt2026protocol}. Its HCP and cellular Hodge primitives
are those of \cite{vanDerVlugt2026hcp}; the projection, reconstruction
and source-distance comparisons are proved here.
The coding theorems do not assume an HCP presentation or a commitment
protocol. A protocol using them must enforce the fixed-pair, later-challenge
order described in Section~6. Section~7 uses Paper 2's explicit systematic
cellular construction and proves its source-distance comparison. It retains
the original committed-cochain payload in a complete coordinate
layout and defines source validity intrinsically by an integral opening and
carry satisfying \(\delta_0a=z+qm\). Scalar cochain edits count \(2n\)
coordinates, while plane and RS edits count \(n\) whole pairs. The proved
global comparison is
\[
  d_C(z)\le d_P(z)=D_2(\bar z)\le2d_C(z).
\]
The factor two is necessary in the near direction; the far implication has
unit scale. An actual well-formed raw payload at
\((q,n,k)=(12289,4096,1024)\) has all three distances equal to \(3/4\).
This is a modular source theorem. The Euclidean and full-HCP residuals have
different comparison obligations, and positive-far raw payloads are distinct
from already accepted exact storage anchors. The selected application field
does not inherit the large-field family's \(2^{-128}\) guarantee.

\section{Problem Setup}

\subsection{Reed--Solomon parameters}

Let \(\F\) be a finite field and let
\(\mathcal L=(a_i)_{i\in I}\) be an injectively indexed evaluation domain,
where \(|I|=n\).  For \(1\le k\le n\), define
\[
  \mathcal C=\RS[\F,\mathcal L,k]
  =
  \left\{
    (f(a_i))_{i\in I}:
    f\in\F[X],\ \deg f<k
  \right\}.
\]
The rate is \(\rho=k/n\).  The smooth-domain regime takes \(\rho\) from a fixed
constant set, such as
\[
  \left\{\frac12,\frac14,\frac18,\frac1{16}\right\},
\]
and requires \(\mathcal L\) to have an appropriate smooth structure.

\subsection{Relative Hamming distance}

For a finite alphabet \(A\) and \(u,v\in A^I\), let
\[
  d_H(u,v)=|\{i\in I:u_i\ne v_i\}|,
  \qquad
  \Delta(u,v)=\frac{d_H(u,v)}{n}.
\]
For a nonempty code \(\mathcal D\subseteq A^I\), define
\[
  \Delta(w,\mathcal D)
  =
  \min_{c\in\mathcal D}\Delta(w,c).
\]
Unless stated otherwise, radii are relative radii in \([0,1]\).
The uniform MCA bound also allows arbitrary real values of its radius
parameter. Natural-number radii describe the finite Hamming grid and its
rounding; distances in the main statements are normalized by \(n\).

\subsection{Interleaving}

For a positive integer \(m\), define
\[
  \mathcal C^{\equiv m}
  =
  \left\{
    p\in(\F^m)^I:
    \text{for every }j\in\{1,\ldots,m\},
    \ (p_{i,j})_{i\in I}\in\mathcal C
  \right\}.
\]
Hamming distance on \((\F^m)^I\) counts a coordinate \(i\) as erroneous when
the complete vector symbol differs.  It does not sum componentwise scalar
distances.

\subsection{List-decoding balls}

For \(w\in(\F^m)^I\) and \(\delta\in[0,1]\), define
\[
  \List(\mathcal C^{\equiv m},w,\delta)
  =
  \left\{
    c\in\mathcal C^{\equiv m}:
    \Delta(w,c)\le\delta
  \right\}.
\]
The worst-case list size is
\[
  L_m(\mathcal C,\delta)
  =
  \max_{w\in(\F^m)^I}
    |\List(\mathcal C^{\equiv m},w,\delta)|.
\]
The maximum over received words is essential.  A bound on a structured subset
of received words is an application result, not the universal list-decoding
statement.

\section{Planes, Pair Words, and Affine Lines}

\subsection{The full ambient pair space}

Define a field-valued two-coordinate plane to be a function
\[
  p:I\longrightarrow(\{0,1\}\longrightarrow\F).
\]
By reassociating function spaces, the plane carrier is exactly
\[
  I\to(\{0,1\}\to\F)
  \cong
  (\F^2)^I.
\]
Thus the word ``plane'' introduces no restriction on the received-word space.

\begin{proposition}[Plane--pair equivalence]
\label{prop:plane-pair-equivalence}
The map sending a plane \(p\) to the vector word
\[
  i\longmapsto(p(i,0),p(i,1))
\]
is a bijection from field-valued planes to \((\F^2)^I\).  It preserves and
reflects Hamming distance.
\end{proposition}

\begin{proof}
The inverse sends a pair word \(w\) to the function
\[
  (i,j)\longmapsto w_i(j).
\]
Both composites are pointwise identities.  At coordinate \(i\), the two planes
are equal exactly when their complete pair symbols are equal, so the Hamming
disagreement sets coincide.
\end{proof}

\subsection{Pair words as affine lines}

For a plane \(p\), define its base and direction by
\[
  p_0(i)=p(i,0),
  \qquad
  p_1(i)=p(i,1).
\]
The associated affine line of scalar received words is
\[
  \ell_p(\gamma)=p_0+\gamma p_1,
  \qquad \gamma\in\F.
\]

\begin{proposition}[Plane--line equivalence]
\label{prop:plane-line-equivalence}
The map \(p\mapsto(p_0,p_1)\) is a bijection between field-valued planes and
affine-line presentations consisting of an arbitrary base word and arbitrary
direction word in \(\F^I\).
\end{proposition}

\begin{proof}
Given an arbitrary pair \((f_0,f_1)\), define
\[
  p(i,0)=f_0(i),\qquad p(i,1)=f_1(i).
\]
This construction is inverse to taking base and direction.
\end{proof}

\subsection{Challenge evaluation is lossy}

For fixed \(\gamma\), the map
\[
  (p_0,p_1)\longmapsto p_0+\gamma p_1
\]
is generally many-to-one. Consequently, distance of the pair from the code
cannot be inferred
deterministically from one fixed challenge without an additional good-challenge
hypothesis.  MCA supplies a probabilistic statement over a challenge sampled
after the pair has been fixed.

\section{The Capacity-Radius Regime}

\subsection{The finite Hamming grid}

Let \(1\le k\le n\).  The capacity endpoint is already an attainable
normalized Hamming radius:
\[
  r_{\mathrm{cap}}=n-k,
  \qquad
  \deltacap=\frac{n-k}{n}=1-\rho.
\]
There is therefore no rounding loss at the endpoint.  For example, the
rate-\(1/4\), length-\(4096\) member has
\[
  k=1024,\qquad r_{\mathrm{cap}}=3072,\qquad
  \deltacap=\frac34.
\]
The endpoint theorem below includes the strictly smaller radius
\(3071/4096\).

There is, however, a distinct granularity issue in the phrase ``largest
relative radius.''  A closed Hamming ball changes only at
\[
  \mathcal H_n=\left\{0,\frac1n,\ldots,\frac nn\right\}.
\]
Consequently, a safe-radius set defined over all real
\(\delta\in[0,1]\) may be half-open and have no largest element.  We report
both the largest radius in \(\mathcal H_n\) and the supremum over real radii.

\subsection{An explicit smooth large-field family}
\label{sec:explicit-family}

For \(u\in\mathbb N\), set
\[
  k_u=64\cdot2^u=2^{u+6},
  \qquad
  n_u=4k_u=256\cdot2^u=2^{u+8},
  \qquad
  q_u=3^{n_u}.
\]
Let \(\F_u=\mathbb F_{q_u}\).  Since the multiplicative group of a finite
field is cyclic of order \(q_u-1\), it suffices to check
\(n_u\mid(q_u-1)\).  Euler's theorem gives this directly:
\[
  \varphi(n_u)=\frac{n_u}{2},
  \qquad
  \gcd(3,n_u)=1,
  \qquad
  3^{n_u/2}\equiv1\pmod{n_u}.
\]
Squaring yields \(3^{n_u}\equiv1\pmod{n_u}\).  Thus
\(\F_u^\times\) contains a subgroup \(H_u\) of order \(n_u\).  We take
\(\mathcal L_u=H_u\) and
\[
  \mathcal C_u=\RS[\F_u,H_u,k_u].
\]
The evaluation length and degree bound are powers of two, the domain is a
multiplicative subgroup, and the rate is exactly \(1/4\).  This is the
standard smooth-domain condition for FFT-based Reed--Solomon codes
\cite{arnonChiesaFenziYogev2024}, and it is the explicit family used in the
quantitative statement.

\subsection{The field-size inequality}
\label{sec:field-inequality}

Fix \(\epsstar=2^{-128}\).  For the family above,
\[
  \binom{n_u}{k_u}\le 2^{n_u}.
\]
Moreover \(n_u\) is even and \(n_u/2\ge128\), so
\[
\begin{aligned}
  q_u=3^{n_u}
  &=(3^2)^{n_u/2}
    =9^{n_u/2} \\
  &\ge8^{n_u/2}
    =2^{3n_u/2}
    \ge2^{n_u+128}.
\end{aligned}
\]
It follows that
\[
  \binom{n_u}{k_u}
  \le 2^{-128}q_u
  =\epsstar|\F_u|.
\]

The pair-MCA theorem uses the comparison
\[
  2^{128}\binom{n_u}{k_u+1}\le q_u.
\]
It follows from the same calculation, because
\(\binom{n_u}{k_u+1}\le2^{n_u}\).
Compared with the ordered certificate count, this removes a factor \(k_u+1\):
\[
  \binom{n_u}{k_u}(n_u-k_u)
  =(k_u+1)\binom{n_u}{k_u+1}.
\]
The field is exponential in the block length.  This is a valid explicit
sufficiently-large-field regime, not a claim about practically sized fields
or about every smooth code supplied to the Grand Challenge.

\section{The Capacity-Grid List Theorem}

The following argument applies to every injectively indexed Reed--Solomon
evaluation domain; smoothness enters only when we instantiate the explicit
family from \cref{sec:explicit-family}.

\begin{theorem}[Universal capacity-grid equality]
\label{thm:capacity-grid-exact}
Let \(\mathcal C=\RS[\F,\mathcal L,k]\) have length \(n\), where
\(1\le k\le n\), and let \(m\ge1\).  Then
\[
  L_m\left(\mathcal C,\frac{n-k}{n}\right)=\binom nk .
\]
In particular, the worst-case endpoint list size is independent of \(m\).
\end{theorem}

\begin{proof}
We first prove the upper bound.  Fix \(w\in(\F^m)^n\).
Every codeword in the ball agrees with \(w\) on at least \(k\) complete
vector-symbol coordinates.  Assign to each such codeword one \(k\)-subset of
its agreement coordinates.  This assignment is injective: if two interleaved
codewords are assigned the same set \(S\), then they agree with one another on
\(S\).  In each component, two degree-\(<k\) polynomials agree at the \(k\)
distinct evaluation points indexed by \(S\), and hence are equal.  The two
interleaved codewords are therefore equal.  There are exactly \(\binom nk\)
possible sets \(S\).

For the reverse inequality, choose one of the \(m\) components and define the
received word
\[
  w_i^\star=(a_i^k,0,\ldots,0).
\]
For \(k<n\), the degree-\(k\) component is a classical RS deep-hole
witness \cite[Section~2]{ketiWan2016deepHoles}. The interpolation argument
below counts its full interleaved endpoint list explicitly.
For each \(k\)-subset \(S\subseteq I\), put
\[
  V_S(X)=\prod_{i\in S}(X-a_i),
  \qquad
  f_S(X)=X^k-V_S(X).
\]
The polynomial \(V_S\) is monic of degree \(k\), so its leading term cancels
that of \(X^k\) and \(\deg f_S<k\).  Use \(f_S\) in the distinguished
component and zero polynomials in all other components to obtain an
interleaved codeword \(c_S\).

At evaluation coordinate \(i\),
\[
  c_{S,i}=w_i^\star
  \quad\Longleftrightarrow\quad
  V_S(a_i)=0
  \quad\Longleftrightarrow\quad
  i\in S,
\]
where the final equivalence uses injectivity of the evaluation points.
Consequently \(c_S\) has distance exactly \(n-k\) from \(w^\star\), and its
agreement set recovers \(S\).  The \(\binom nk\) choices of \(S\) therefore
give distinct codewords in the endpoint ball.  This matches the upper bound.
\end{proof}

\begin{theorem}[Failure on the next Hamming shell]
\label{thm:next-shell-lower}
Under the same hypotheses,
\[
  \left|
    \List\left(
      \mathcal C^{\equiv m},
      0,
      \frac{n-k+1}{n}
    \right)
  \right|
  \ge |\F|^m .
\]
\end{theorem}

\begin{proof}
Fix \(k-1\) evaluation points \(S\subseteq\mathcal L\) and set
\[
  g_S(X)=\prod_{a\in S}(X-a).
\]
For every \(\lambda=(\lambda_1,\ldots,\lambda_m)\in\F^m\), the tuple
\[
  (\lambda_1g_S,\ldots,\lambda_mg_S)
\]
defines an interleaved codeword.  These \(|\F|^m\) codewords are distinct.
The zero tuple has distance zero from the received word \(0\); every nonzero
tuple vanishes on \(S\) and is nonzero at every evaluation point outside
\(S\), so its distance is exactly \(n-k+1\).
\end{proof}

\begin{theorem}[Fractional-target phase split]
\label{thm:fractional-phase-split}
Fix \(0<\epsstar<1\), write \(q=|\F|\), and call a radius \(\delta\)
safe when
\[
  L_m(\mathcal C,\delta)\le\epsstar q.
\]
For every \(m\ge1\), exactly one of the following regimes applies:
\begin{enumerate}[label=(\roman*)]
  \item If \(\epsstar q<1\), no radius in \([0,1]\) is safe.
  \item If
    \(1\le\epsstar q<\binom nk\), radius zero is safe but the capacity
    endpoint and every larger radius fail.  The optimal sub-endpoint radius is
    not determined by the present argument.
  \item If \(\binom nk\le\epsstar q\), the largest safe radius on
    \(\mathcal H_n\) is
    \[
      \deltacap=\frac{n-k}{n}=1-\rho;
    \]
    and over real radii the safe set is exactly
    \[
      \left[0,\frac{n-k+1}{n}\right),
    \]
    so it has no largest element and has supremum
    \((n-k+1)/n\).
 \end{enumerate}
Equivalently, the fractional target holds at the capacity endpoint if and
only if
\[
  \binom nk\le\epsstar|\F|.
\]
\end{theorem}

\begin{proof}
Every codeword is contained in the radius-zero ball centered at itself, so no
radius can be safe when \(\epsstar q<1\).  A radius-zero ball contains at most
one codeword, proving its safety when \(\epsstar q\ge1\).

\Cref{thm:capacity-grid-exact} proves both the failure in the intermediate
regime and the necessary-and-sufficient endpoint criterion.  In the final
regime, it and monotonicity of list balls prove safety through count radius
\(n-k\).  A real radius strictly below
\((n-k+1)/n\) defines the same ball as some count radius at most \(n-k\).
At the next grid point, \cref{thm:next-shell-lower} gives a list of size at
least \(|\F|^m\ge|\F|>\epsstar|\F|\); monotonicity then rules out every
larger radius.
\end{proof}

\begin{corollary}[The explicit smooth family]
\label{cor:explicit-family-list}
For every \(u\ge0\), every \(m\ge1\), and
\(\mathcal C_u=\RS[\F_u,H_u,k_u]\) from
\cref{sec:explicit-family}, the conclusions of
\cref{thm:fractional-phase-split}\textup{(iii)} hold with
\(\epsstar=2^{-128}\).  In
particular, the largest safe Hamming-grid radius is \(3/4\).
\end{corollary}

\begin{proof}
Apply the inequality proved in \cref{sec:field-inequality}.
\end{proof}

\subsection{Scope relative to the Grand Challenge}

\Cref{thm:capacity-grid-exact} completely determines the endpoint list size
for every supplied Reed--Solomon code, smooth or otherwise.  Thus the displayed
binomial field-size inequality is necessary and sufficient at that endpoint,
not merely a convenient sufficient hypothesis.
\Cref{cor:explicit-family-list} supplies a concrete smooth family in the safe
regime.  The present argument does not determine the largest sub-endpoint
radius for an arbitrary fixed smooth code in the intermediate regime
\[
  1\le2^{-128}|\F|<\binom nk.
\]
The exponential alphabet is what places the explicit family in the final
regime; it is not needed for the universal endpoint equality.

This distinction matters in light of known capacity results.  Generic
Reed--Solomon evaluation points can, for every fixed positive gap
\(\epsilon\), achieve radius \(1-\rho-\epsilon\) over exponentially large
fields with a list size bounded in terms of \(\epsilon\)
\cite{brakensiekGopiMakam2024}. This positive-gap statement differs from the
exact endpoint equality above, and generic points do not supply the
structured smooth domain required here. Conversely, near-capacity proximity-gap
behavior can fail for particular smooth prime-field families
\cite{kambire2026}.  The theorem above imports neither result: it uses only
polynomial interpolation and an explicit field-size comparison.

Recent work of Dao, Kominers, and Thaler \cite{daoKominersThaler2026}
gives quantitative list-decoding and MCA bounds on prescribed evaluation
domains, including beyond-Johnson bounds under sufficiently large
characteristic. Their positive agreement-gap counting results do not cover
the exact agreement endpoint \(A=k\); their beyond-Johnson counting
hypotheses also exclude the characteristic-three family with \(k\ge64\)
used here. Their Johnson-regime results in every characteristic remain a
separate comparison. The present binomial bounds are not asserted to improve
the sharper radius-dependent estimates in that literature, which also
includes existing Lean/ArkLib formalizations.

\section{Pair Mutual Correlated Agreement}

\subsection{Bad challenges}

Let \(p=(p_0,p_1)\in(\F^2)^I\), and write
\[
  y_\gamma=p_0+\gamma p_1.
\]
For \(S\subseteq I\), say that \(S\) is \(\delta\)-large when
\[
  \frac{|S|}{n}\ge1-\delta.
\]
Define \(\Bad_{\MCA}(\mathcal C,\delta,p)\) to be the set of
\(\gamma\in\F\) for which there exists a \(\delta\)-large set \(S\) such that:
\begin{enumerate}[label=(\roman*)]
  \item \(y_\gamma\) agrees on \(S\) with a codeword of \(\mathcal C\);
  \item the pair \(p\) does not agree on the same \(S\) with a codeword of
    \(\Cpair\).
\end{enumerate}
The MCA error is
\[
  \varepsilon_{\MCA}(\mathcal C,\delta)
  =
  \max_{p\in(\F^2)^I}
  \frac{|\Bad_{\MCA}(\mathcal C,\delta,p)|}{|\F|}.
\]

This is the two-word same-set event appearing in
\cite[Theorem~4.6, full version]{benSassonCarmonHabockKoppartySaraf2026}.
That source uses polynomials of degree at most its parameter, so its
parameter is \(k-1\) in the degree-\(<k\) convention used here.
Our bound below is uniform in the radius; the cited literature gives
sharper radius-dependent bounds in its stated regimes. At fixed rate and a
fixed positive gap below the Johnson radius, the cited theorem gives an
exceptional-challenge count linear in \(n\), generally much smaller than
the binomial bound below. Broader polynomial-generator formulations of MCA
for linear codes are studied by Bordage et al.
\cite{bordageEtAl2026polynomialMCA}.

\begin{theorem}[Uniform pair-MCA certificate bound]
\label{thm:pair-mca}
Let \(\mathcal C=\RS[\F,\mathcal L,k]\) have injectively indexed
evaluation domain of length \(n\), where \(1\le k\le n\).  For every
\(\delta\in\mathbb R\) and every \(p\in(\F^2)^n\),
\[
  |\Bad_{\MCA}(\mathcal C,\delta,p)|
  \le \binom n{k+1}.
\]
Consequently,
\[
  \varepsilon_{\MCA}(\mathcal C,\delta)
  \le \frac{\binom n{k+1}}{|\F|}.
\]
Smoothness is not required for this bound.
\end{theorem}

\begin{proof}
Fix a bad challenge \(\gamma\), and choose a witnessing set \(S\).
Necessarily \(|S|>k\).  Indeed, when \(|S|\le k\), interpolate each component
of \(p\) on \(S\) by a polynomial of degree \(<k\); the resulting pair
codeword agrees with \(p\) on \(S\), contradicting condition~(ii).

Choose a \(k\)-subset \(T\subset S\).  Let \(a_T,b_T\) be the unique
degree-\(<k\) polynomials interpolating \(p_0,p_1\), respectively, on \(T\).
Condition~(ii) says that the pair interpolant \((a_T,b_T)\) fails somewhere
on \(S\). Choose such an \(i\in S\setminus T\), and set
\(U=T\cup\{i\}\). Condition~(i) remains true after restriction from \(S\) to
\(U\). The pair still cannot agree with \(\Cpair\) on \(U\): any purported
pair codeword would agree with \((a_T,b_T)\) on \(T\), hence equal it by
\(k\)-point determination, contradicting the mismatch at \(i\).
Thus \(U\) is a \((k+1)\)-coordinate obstruction for \(\gamma\).

A fixed \(U\) obstructs at most one challenge. Suppose distinct
\(\gamma,\eta\) were both bad on \(U\), with scalar codewords
\(c_\gamma,c_\eta\) agreeing with
\(p_0+\gamma p_1,p_0+\eta p_1\), respectively. By linearity,
\[
  d=\frac{c_\gamma-c_\eta}{\gamma-\eta},
  \qquad
  a=c_\gamma-\gamma d
\]
are codewords. On \(U\), they agree with \(p_1\) and \(p_0\), respectively,
so \((a,d)\) is a pair codeword agreeing with \(p\) on \(U\), a
contradiction.

Make the finite choices deterministic. Mapping each bad challenge to its
unordered obstruction \(U\) is injective, and there are
\(\binom n{k+1}\) such subsets.
Dividing by \(|\F|\) proves the error bound.  Notice that the large-set
inequality is no longer used after the witness is extracted; this is why the
bound is uniform in \(\delta\).
\end{proof}

\begin{corollary}[Explicit smooth family at \(2^{-128}\)]
\label{cor:explicit-family-mca}
For every member \(\mathcal C_u\) of the family in
\cref{sec:explicit-family}, every \(\delta\in[0,1]\), and every pair word
\(p\),
\[
  \frac{|\Bad_{\MCA}(\mathcal C_u,\delta,p)|}{|\F_u|}
  \le 2^{-128}.
\]
Thus, for this family and this same-set pair-MCA definition, every normalized
radius is safe and the largest admissible radius is \(1\).
\end{corollary}

\begin{proof}
\Cref{sec:field-inequality} proves
\[
  2^{128}\binom{n_u}{k_u+1}\le q_u.
\]
Apply \cref{thm:pair-mca}.
\end{proof}

\subsection{An exact small-radius boundary instance}
\label{sec:mca-boundary-field}

The uniform certificate bound is sufficient but need not locate the first
unsafe radius.  Write
\[
  B_{\mathcal C}(r)
  =
  \max_p|\Bad_{\MCA}(\mathcal C,r/n,p)|
\]
for the exact worst bad-challenge count at count radius \(r\).  Consider
\[
  q=2^{128}+192513
   =340282366920938463463374607431768403969.
\]
The accompanying Lean development certifies that \(q\) is prime and
\(q\equiv1\pmod {4096}\).  Hence \(\mathbb F_q^\times\) contains an
order-\(4096\) subgroup, which we use as the evaluation domain for the smooth
RS code with \(n=4096\) and \(k=1024\).

\begin{theorem}[Exact 128-bit small-radius transition]
\label{thm:mca-boundary-threshold}
For this code,
\[
  B_{\mathcal C}(0)=1,
  \qquad
  B_{\mathcal C}(1)\ge2.
\]
Consequently, for the \(2^{-128}\) MCA error target, its greatest safe count
radius is \(0\) and its first unsafe count radius is \(1\).
For the relative-radius definition over
all \(\delta\in[0,1]\), the safe set is
\[
  \left[0,\frac1{4096}\right),
\]
so the real-radius supremum is \(1/4096\) and is not attained.
\end{theorem}

\begin{proof}
At count radius zero, every large agreement set is the whole coordinate set.
A fixed set supports at most one bad challenge, while a unit-spike affine
line realizes one; hence \(B_{\mathcal C}(0)=1\).  At count radius one, take
two coordinates outside a \(k\)-point interpolation set.  The affine line
from the unit spike at the first coordinate to the unit spike at the second
has distinct bad challenges \(0\) and \(1\), proving
\(B_{\mathcal C}(1)\ge2\).

The certified field lies in the exact budget window
\[
  2^{128}\le q<2^{128}\cdot2.
\]
Thus the integer criterion \(2^{128}B_{\mathcal C}(r)\le q\) holds at
\(r=0\) and fails at \(r=1\).  Monotonicity of the bad-count spectrum gives
the two grid endpoints.

Finally, if \(0\le\delta<1/n\), the inequality
\(|S|/n\ge1-\delta>1-1/n\) forces the integer \(|S|\) to equal \(n\).
The bad sets are therefore exactly those at count radius zero throughout
that half-open interval.  At \(\delta=1/n\), sets missing one coordinate
become admissible, and the two-spike line already violates the budget.
Monotonicity preserves failure for larger radii, proving the real-radius
claim.
\end{proof}

\subsection{Distance interpretation}

Let
\[
  D_2(p)=\Delta(p,\Cpair),
  \qquad
  D_1(y)=\Delta(y,\mathcal C).
\]
For a challenge outside the MCA bad set, scalar closeness on a large agreement
set forces pair closeness on that same set. For \(\delta\in[0,1]\), this
gives the distance implication
\[
  D_2(p)\ge\delta
  \Longrightarrow
  D_1(y_\gamma)\ge\delta
\]
for every MCA-good \(\gamma\), with the fraction of excluded challenges bounded
by \(\varepsilon_{\MCA}(\mathcal C,\delta)\).

To see the boundary convention explicitly, suppose that
\(D_1(y_\gamma)<\delta\), and choose a closest scalar codeword \(c\).
Let \(S=\{i:y_\gamma(i)=c(i)\}\). Then
\(|S|/n=1-D_1(y_\gamma)>1-\delta\), so \(S\) is \(\delta\)-large.
Since \(\gamma\) is MCA-good, a pair codeword agrees with \(p\) on
\(S\). Hence \(D_2(p)\le D_1(y_\gamma)<\delta\). Taking the
contrapositive proves the displayed implication, including its non-strict
\(\ge\) boundary.

This is a challenge-indexed statement.  Replacing the sampled challenge by one
fixed field element would require a separate theorem that the selected element
is good for every pair, which does not follow from a fractional MCA bound.

\subsection{Endpoint interpretation}

At \(\deltacap=1-k/n\), interpolation on any \(k\) coordinates shows that
\[
  D_1(y)\le\deltacap
  \quad\text{and}\quad
  D_2(p)\le\deltacap
\]
for every scalar word \(y\) and pair word \(p\).  Moreover, on every set
\(S\) of size at most \(k\), each component of \(p\) extends to a
degree-\(<k\) polynomial.  Such an \(S\) can therefore never witness an MCA
failure.  Every bad endpoint challenge must instead have a witness with
\(|S|\ge k+1\).

Consequently, the endpoint distance implication above is an equality
non-collapse statement: when \(D_2(p)=\deltacap\), a good challenge must not
lower the projected distance strictly below \(\deltacap\).
\Cref{thm:pair-mca} supplies a uniform quantitative bound on the
excluded challenges.

\subsection{Commit-then-challenge order}

The mathematical experiment fixes \(p\) before sampling
\(\gamma\leftarrow\F\).  This ordering prevents the pair from being selected as
a function of the challenge.  A full protocol implementation must bind \(p\)
before deriving \(\gamma\), but the coding theorem needs only the resulting
quantification and probability space; it does not depend on a complete
blockchain or commitment protocol.

\subsection{Scope}

The theorem is universal over pair words and injective RS evaluation domains,
but its \(2^{-128}\) corollary uses the explicit exponential-size field.
For a different supplied field, the same certificate proof gives the
sufficient comparison
\[
  \binom n{k+1}\le 2^{-128}|\F|.
\]
Relative to the older ordered-certificate count used in this development,
the bound saves the exact factor \(k+1\), since
\(\binom nk(n-k)=(k+1)\binom n{k+1}\).
This is not a uniform improvement over the cited literature bounds.
If that inequality fails, this proof does not determine the optimal radius or
show that the MCA target is impossible; it only ceases to certify the requested
error.

For an arbitrary supplied finite code, radius \(r/n\) meets the
\(2^{-128}\) target exactly when
\[
  2^{128}B_{\mathcal C}(r)\le|\F|.
\]
This is also the exact formal equivalence: the Lean theorem identifies the
real-valued operational error predicate at normalized radius \(r/n\) with the
displayed integer inequality, so no rounding relaxation is hidden here.
Thus a uniform upper-bound certificate at \(r\), together with one concrete
affine line exceeding the budget at \(r+1\), pins both the largest safe and
first unsafe Hamming-grid radii. The remaining middle regime is the
code-specific problem of proving those matching bounds; smoothness by itself
does not determine this incidence spectrum.

The result is for the two-coordinate, same-set MCA experiment defined above.
Broader notions or higher-dimensional challenge families require
separate statements.

As a separate concrete/application result, the repository also checks the
canonical runtime field \(q=12289\) with \(n=4096\) and \(k=1024\).  Since
\(q<2^{128}\) while its radius-zero unit-spike line has one bad challenge,
that concrete code has no safe count radius for the \(2^{-128}\) MCA error
target. This
runtime-facing fact is not used by the standalone large-field or boundary-field
theorems above.

\section{Application Interface}

Paper 2 specifies the shared cochain construction, snapshot/decoder context,
coordinate layout and public observation profile
\cite{vanDerVlugt2026protocol}. This section recalls only the data needed to
state the projection and proves its quantitative validity, distance and
challenge comparisons. The commitment-security, protocol gate/evidence and
state-transition contracts remain with that construction. Paper 1 supplies
the HCP and cellular Hodge primitives \cite{vanDerVlugt2026hcp}; the
exact-component calibration obstruction is proved in
\cref{prop:exact-component-obstruction} below.

The coding theorem is stated over the complete received-word space and does not
depend on a particular source of observations.  An application may nevertheless
inherit the result through a distance-aware map.

Let \(O\) be a type of raw presented observations, let
\(G\subseteq O\) be the exact/valid locus, and let
\[
  \pi:O\longrightarrow(\F^2)^I
\]
be the extracted pair word.  A source distance
\(\rho_O:O\to[0,\infty)\) is meaningful only when \(O\) includes arbitrary or noisy
presentations; if every element of \(O\) already contains a validity witness,
then the proximity problem may collapse to the exact code locus.

The Lean companion packages this interface as a source validity model and a
source-to-code calibration.  Its only quantitative field is the pointwise
inequality
\[
  \rho_O(o)\le D_2(\pi(o)).
\]
It yields, uniformly in \(\delta\),
\[
  \rho_O(o)\ge\delta
  \Longrightarrow
  D_2(\pi(o))\ge\delta.
\]

Under this comparison, the proved universal list theorem applies to
\(\pi(o)\); the proved pair-MCA theorem applies through the same interface
once the application fixes its pair before sampling the challenge.

There are two standard ways to discharge the comparison. Defining
\(\rho_O=D_2\circ\pi\) makes it reflexive. A pre-existing source
metric requires a mathematical proof of
\(\rho_O\le D_2\circ\pi\).  Exact witnesses may certify membership in \(G\),
but the carrier \(O\) must still include arbitrary received observations.

Protocol-specific observation formats, commitments, and runtime verification
are intentionally outside the proof of
\cref{thm:capacity-grid-exact,thm:next-shell-lower,thm:pair-mca}.
Such developments may cite the two coding theorems through the interface
above; neither theorem assumes that every received word comes from an
application-specific witness space.

\subsection{Projective coordinates in the projection chain}\label{sec:projective-observation-boundary}

Let \(V=\mathbb R^n\) have its fixed Euclidean inner product, and let
\(w_1,\ldots,w_r\) be an independent frame spanning \(W\le V\). Its exterior
coordinate is
\[
  c=w_1\wedge\cdots\wedge w_r\in\bigwedge^r V.
\]
An independent frame has nonzero exterior product. A different basis of the
same subspace changes \(c\) by the nonzero determinant of its change-of-basis
matrix. The projective class therefore depends only on \(W\):
\[
 c\sim c'\quad\Longleftrightarrow\quad
 \exists u\in\mathbb R^\times,\ c'=uc.
\]
The raw carrier is the scalar-orbit quotient of the whole exterior power,
including its zero orbit. The conclusions below use the nonzero decomposable
orbits determined by genuine frames.

A related specialized reconstruction uses restricted Pl\"ucker coordinates
for maximal isotropic subspaces \cite{dingKim2026plucker}. The argument here
uses the full projective point of an arbitrary real subspace at the fixed
ambient metric.

\begin{proposition}[Projective Pl\"ucker invariance]\label{prop:plucker-invariance}
  Two independent frames spanning the same finite-dimensional subspace have
  equal projective exterior coordinates.
\end{proposition}
\begin{proof}
  Both exterior products are nonzero elements of the one-dimensional top
  exterior power of that subspace. They differ by a nonzero scalar and hence
  have the same orbit.
\end{proof}

\paragraph{Recovery from the full point.}

\begin{proposition}[Recovery from the full projective point]
\label{prop:plucker-recovery}
  Let \(w_1,\ldots,w_r\) be an independent frame for
  \(W\subseteq\mathbb R^n\), and put
  \(c=w_1\wedge\cdots\wedge w_r\). Then
  \[
    W=\{x\in\mathbb R^n:x\wedge c=0\}.
  \]
  This reconstruction depends only on \([c]\). Consequently, two genuine
  rank-\(r\) frame witnesses have the same full projective point if and only
  if they present the same subspace. At the fixed ambient inner product,
  the full point determines \(P_W\) and \(I-P_W\).
\end{proposition}
\begin{proof}
  If \(x\in W\), the frame with \(x\) adjoined is dependent and its exterior
  product vanishes. If \(x\notin W\), the adjoined frame is independent and
  its exterior product is nonzero. Multiplying \(c\) by a nonzero scalar
  does not change this annihilator. The projector is the unique orthogonal
  projection onto the reconstructed subspace.
\end{proof}

The zero orbit has the whole ambient space as its annihilator, which is why
the statement requires a genuine frame. For the operative harmonic subspace
\(W=\ker\Delta_x\), the reconstructed projector and \(b\) give
\[
  h=P_Wb.
\]
The positive spectrum of \(\Delta_x\) remains additional operator data.
Replacing the full real coordinates by encoded data or a digest requires a
separate reconstruction procedure.

The geometric source can enter the existing extraction chain through
\[
  [w_1\wedge\cdots\wedge w_r]
  \xrightarrow{\operatorname{enc}} c
  \xrightarrow{E} p_c
  \xrightarrow{\bmod q} \bar p_c
  \xrightarrow{\gamma} y_\gamma.
\]
Here \(\operatorname{enc}\) is the supplied projective-to-cochain map,
\(c\) is the original commitment payload, and the remaining maps are the
actual plane extraction, coefficient reduction and challenge evaluation
specified below. The first map requires its own realization and validity
hypotheses; reconstruction of a subspace does not construct that codec.
The projection-chain theorem composes the supplied validity and near-distance
comparisons by applying each stage in order and using transitivity of the
distance bounds. A source-far conclusion requires a separate comparison in
the direction used by the application calibration; it is not manufactured
by composing these near-side maps. Full-coordinate reconstruction supplies
no such quantitative inequality.

\subsection{The existing committed-cochain projection}

The project-specific application uses Paper 2's existing commitment payload
\cite{vanDerVlugt2026protocol},
\[
  c=(q,N,\mathbf c,\tau).
\]
Here \(\mathbf c\) is the serialized integer edge-value list. The Lean
carrier \lean{CommittedCochain} is an abbreviation for the same
\lean{CochainCommitmentData} used by the commitment and artifact
layers. A selected observation exposes that payload; it does not replace it
with a second commitment object. The construction and authentication
conditions of that observation belong to Paper 2; here the payload is the
input to the actual extraction map.

Fix an extraction context consisting of an expected edge count, selected
indices \(\iota(i,j)\), and a coefficient interpretation
\(a:\mathbb Z\to\mathbb Z\). For an aligned payload, the actual plane
projection is
\[
  p_c(i,j)=a(\mathbf c_{\iota(i,j)}),
  \qquad j\in\{0,1\}.
\]
Alignment ensures that every selected index is in bounds. The public
plane-view certificate records equality with this extraction from the
observed payload. The canonical instantiation reduces both coordinates
modulo \(q=12289\), with RS length \(n=4096\) and degree bound \(k=1024\):
\[
  \bar p_c(i,j)=p_c(i,j)\bmod q,
  \qquad
  y_\gamma(i)=\bar p_c(i,0)+\gamma\bar p_c(i,1).
\]
Thus one observed commitment supplies the full pair before the challenge is
evaluated. The \lean{PublicCalibration} module implements this extraction
and proves its equality with the observed-payload projection.

On integer planes whose coordinates lie in \([0,q)\), coordinatewise
reduction is a bijection onto the field-pair space. It preserves whole-pair
Hamming distance. Let \(\mathcal C_A\) be the pair code pulled back to
this canonical integer alphabet, and let \(D_A\) be relative Hamming
distance to \(\mathcal C_A\). The existing Lean development proves
\[
  D_A(p_c)=D_2(\bar p_c).
\]
The public-observation calibration uses \(D_A(p_c)\) as its source score,
so its all-radius farness implication is already proved. Its observation
carrier requires aligned, canonical-region payload views without requiring
pair-codeword membership.

An application retaining an independently chosen cochain, Hodge, or
Pl\"ucker disagreement \(\rho_c\) still needs a quantitative comparison,
for example
\[
  \rho_c(c)\ge r_S
  \quad\Longrightarrow\quad
  D_A(p_c)\ge r_T.
\]
The equality above does not prove this upstream comparison. In particular,
changing unselected edge values leaves \(p_c\) unchanged; any source metric
sensitive to those values needs an admissible-family restriction or a
separate control theorem.

The storage interface now fixes the same snapshot, payload interpretation
and public artifact in one context, while keeping exact acceptance separate
from raw context binding. For its decoded cochain, one independent geometric
score is
\[
  \rho_{\mathrm{ex}}(c)
   =\|P_{(\image\delta_0)^\perp}\operatorname{dec}_S(c)\|.
\]
It vanishes exactly on accepted exact payloads once the context and shape
conditions hold. A sufficient comparison criterion is an exact comparison cochain
\(\lambda(z)\in\image\delta_0\) for every target pair-codeword \(z\),
together with a constant \(A>0\) and the bound
\[
  \|\operatorname{dec}_S(c)-\lambda(z)\|
    \le A\,\delta_H(z,\bar p_c).
\]
Choosing an attained nearest codeword and applying orthogonal-projection
contractivity gives
\[
  \rho_{\mathrm{ex}}(c)\le A D_2(\bar p_c),
  \qquad
  \rho_{\mathrm{ex}}(c)\ge r_S
    \Longrightarrow D_2(\bar p_c)\ge r_S/A.
\]
The storage calibration module proves this implication from the lift and
decoder bounds. It does not construct those application-specific inputs.
The canonical storage adapter identifies the real-infimum pair distance
with the attained generic minimum and joins both comparisons to the actual
observed artifact.

\subsection{Source-metric obstructions and a no-wrap profile}

Fix a real cellular cochain complex
\(C^0\xrightarrow{\delta_0}C^1\xrightarrow{\delta_1}C^2\)
with its chosen inner products and harmonic subspace
\(\mathcal H=\ker\Delta_1\). The full Hodge sample
\(b=h+\delta_0a+\delta_1^*v\) and its exact-component view
\(\delta_0a\) retain different data.

\begin{proposition}[Exact-component obstruction]
\label{prop:exact-component-obstruction}
  Suppose \(v\in C^2\) and \(\delta_1^*v\ne0\). On the unrestricted
  Hodge-noise carrier, the samples
  \[
    b_t=t\delta_1^*v,\qquad t\in\mathbb R,
  \]
  have zero harmonic message and zero exact component, while
  \[
    \|P_{\mathcal H^\perp}b_t\|
      =|t|\,\|\delta_1^*v\|.
  \]
  Thus this full-observation residual does not factor through the exact
  component. If an extraction depends only on that component and its zero
  image has target distance zero, it cannot preserve any positive target
  far radius uniformly on this unrestricted carrier, at any source radius.
\end{proposition}
\begin{proof}
  A coexact vector is orthogonal to \(\mathcal H\), so its orthogonal
  projection onto \(\mathcal H^\perp\) is itself. Homogeneity gives the
  displayed norm. The exact-component view is constantly zero, whereas
  this norm is zero at \(t=0\) and positive at \(t=1\), proving
  nonfactorization. For any fixed source radius, choose positive \(t\)
  large enough to exceed it. The extracted value still has distance zero,
  contradicting every positive target far radius.
\end{proof}

A restricted noise or realization profile may exclude this family, so a
comparison for such a profile requires its explicit hypotheses. The scores
\[
  \|P_{H^\perp}b\|,\qquad
  \operatorname{dist}(b,\image\delta_0),\qquad
  \text{a modular cochain or code distance}
\]
measure harmonic residual, distance to exactness, and modular or code
disagreement, respectively. A modular distance needs its own carrier and
normalization.

The canonical systematic profile below has no faces and satisfies
\(\delta_1=0\), so the proposition's nonzero-coexact premise is
unavailable in that profile.


Canonical representatives alone do not make a finite-field RS code into a
proper real-linear exact-cochain locus. Suppose \(q\) is prime, the evaluation
points are distinct and \(k\ge2\). For a scalar field word \(w\), write
\(\ell(w)\in\mathbb R^I\) for its canonical coordinate lift. The lifts of
the affine RS words already span \(\mathbb R^I\). Indeed, for each coordinate
\(i\), put \(v_j=a_j-a_i\). Canonical representatives satisfy
\[
 \mathbf1-\ell(v)+\ell(v-\mathbf1)=q e_i.
\]
Only the \(i\)-th coordinate wraps at zero. Both affine words and the
constant word belong to RS, so every coordinate vector belongs to the real
span of the lifted codewords. Applying this argument separately to both
columns proves the corresponding assertion in \(\mathbb R^{I\times\{0,1\}}\).

Consequently, if a real-linear decoder \(T\) has any nonexact output, there
is no constant \(A\) for which
\[
 \operatorname{dist}(T\ell(w),\image\delta_0)\le A D_2(w)
 \quad\text{for every field-pair word }w.
\]
Such a bound would make all lifted pair-codewords exact, hence make the
decoder's entire image exact. This obstruction concerns a global comparison
with no additive loss. It leaves restricted profiles, additive defects and
separately justified threshold implications open.

One explicit restricted profile retains the original Euclidean source metric.
Choose an information set \(S\) of size \(k\) and let \(A_{i,s}\in\mathbb F_q\)
be the actual systematic RS interpolation coefficients. Lift them to integers
\(\widetilde A_{i,s}\in[0,q)\), and define the real generator and integer check
by
\[
 (G_{\mathbb R}u)_i=\sum_{s\in S}\widetilde A_{i,s}u_s,
 \qquad (Hz)_i=z_i-\sum_{s\in S}\widetilde A_{i,s}z_s.
\]
In Paper 2's systematic construction, both columns use the complete coordinate
layout with identity coefficient map. Its two-vertex, one-edge cellular sheaf realizes
\(\image\delta_0=\image(G_{\mathbb R}\oplus G_{\mathbb R})\); this equality is
proved from its restriction maps. The decoder reads the same original
integer payload into its real edge coordinates.

Assume \(|z_{i,j}|\le B\), \(B>0\), and the integer-syndrome condition
\(|(Hz_{\cdot,j})_i|<q\) for all \(i,j\). If the reduced pair is a codeword,
every integer check is divisible by \(q\); the syndrome bound forces it to
vanish. Thus codeword validity reflects geometric exactness. If the source is
nonexact, its extracted pair has distance at least \(1/n\). Since
\(\rho_{\mathrm{ex}}(z)\le\|z\|\le B\sqrt{2n}\), these facts give
\[
 \rho_{\mathrm{ex}}(z)\le Bn\sqrt{2n}\,D_2(\bar z),
 \qquad
 \rho_{\mathrm{ex}}(z)\ge r
 \Longrightarrow D_2(\bar z)\ge\frac{r}{Bn\sqrt{2n}}.
\]
The source score remains Euclidean distance to exact cochains. The no-wrap
condition bounds the integer syndrome, not just individual representatives.
An outside-information-coordinate unit error in one column has nonzero real
syndrome, satisfies this profile and gives a positive source distance.

This construction is a coarse positive-radius example. For a nonempty source
far region its certified target radius is at most \(1/n\); it does not supply
a capacity- or Johnson-scale geometric application theorem. Positive decoded
farness concerns raw received payloads and excludes already accepted exact
anchors. The constructed systematic sheaf is an explicit application model,
not an automatic identification with an arbitrary HCP presentation, and no
encryption-security claim follows from this construction.

\subsection{Three native metrics for modular cochains}

The same shared cellular construction has a global comparison in its native
modular metrics. The generic modular result uses a prime field \(\mathbb F_q\), an
injective RS domain with \(0<k\le n\), and a chosen information set of size
\(k\); subgroup smoothness is not required. Recall that the shared cochain
valid locus is specified intrinsically by an actual integral opening and
integer carry:
\[
 \delta_0a=z+qm.
\]
This is modular exactness, distinct from real decoded exactness. Reducing the
actual integer interpolation coefficients proves that modular exactness is
equivalent to the original RS pair-code predicate. No no-wrap condition is
needed for this modular statement.

Let \(d_C(z)\) be scalar-coordinate edit distance to these modularly exact
cochains, normalized by \(2n\). Let \(d_P(z)\) be whole-pair edit distance
to modularly exact integer planes, normalized by \(n\). The RS distance is
\(D_2(\bar z)\). For two words, their mismatch counts satisfy
\(h_P\le h_C\le2h_P\). Taking attained nearest codewords gives
\[
 d_C(z)\le d_P(z)\le2d_C(z).
\]
The plane/RS distance equality is also a theorem. Reduction cannot create a
mismatch. Conversely, lift a field codeword pair by retaining each already
matching received integer pair, and using canonical representatives at the
remaining positions. The lifted word is modularly exact and has exactly the
field word's mismatched positions. Thus
\[
 \boxed{\ d_C(z)\le d_P(z)=D_2(\bar z)\le2d_C(z)\ }.
\]
The scalar representative-lifting argument also identifies \(d_C\) with
the corresponding flattened finite-field code distance. These independently
defined native distances are compared through the actual maps. A unit-scale
near contract must use the explicitly scaled local score \(2d_C\); the
far implication uses \(d_C\le D_2\).

The factor two is sharp for proper RS codes. Put \(u_i=a_i^k\), and take
the pair \((u,0)\). Its whole-pair distance is \((n-k)/n\), while its
scalar-coordinate distance, normalized by \(2n\), is \((n-k)/(2n)\).
The same representative-lifting equalities transfer this witness to the
native integer metrics, so the upper bound cannot be replaced by a
unit-factor near comparison.

There is a nonempty capacity-scale family in these modular metrics. For
\(0<k<n\), take \(u_i=a_i^k\) and both pair columns equal to \(u\).
A polynomial of degree below \(k\) agrees with \(u\) on at most \(k\)
positions, while interpolation supplies a codeword agreeing on any chosen
\(k\) positions. Applying this separately to both columns, and to their
common interpolation, proves
\[
 d_C(z)=d_P(z)=D_2(\bar z)=\frac{n-k}{n}.
\]
For the canonical instance \((q,n,k)=(12289,4096,1024)\), representatives
of these two columns serialize into the original commitment payload in the
declared complete layout, and all three distances are \(3/4\).
The payload is a well-formed raw proximity input. Under the matching decoder,
accepted exact anchors have zero native distances. This is a native modular source
statement, not a capacity-scale Euclidean Hodge-distance theorem.

For this diagonal pair, the challenged word is \((1+\gamma)u\). Every
\(\gamma\ne-1\) preserves the scalar capacity distance and is MCA-good;
\(\gamma=-1\) collapses to zero. At every nonnegative tested radius its bad
set is exactly \(\{-1\}\), with uniform fraction \(1/q\). The existing
canonical mix is noncancelling because its order is \(4096>2\).

The restriction to the stated line profile is essential. For an arbitrary
fixed nonzero \(\gamma_0\), the pair
\(\bigl(u,-\gamma_0^{-1}u\bigr)\) has scalar-coordinate and pair distance
\((n-k)/n\), but its challenge at \(\gamma_0\) is zero. A global fixed-mix
far-preservation claim therefore fails even on capacity-far pairs. The
general protocol must bind the pair before sampling its challenge or prove
goodness for its actual restricted line profile.

For the full HCP residual,
\cref{prop:exact-component-obstruction} retains the nonzero-coexact and
unrestricted-source hypotheses. An exact-component-only view needs a
justified quotient or restricted source profile before it can support a
full-HCP far-region claim. Applying the proposition to a serializer also
requires proving that its extraction factors through that decoded view.

For scalar transfer, the pair must remain fixed before a random challenge is
sampled, or a chosen fixed challenge must be proved MCA-good. The canonical
fixed mix \(g=1331\) is a different assertion from a bound on the fraction
of bad random challenges. Restricting integer representatives makes the
pair faithful; it does not make one scalar evaluation injective.

The current canonical field is also distinct from the sufficiently large
family of Section~4. For the full canonical pair space, the existing
small-radius certificate already exceeds the \(2^{-128}\) budget at count
radius zero because \(12289<2^{128}\). This does not determine the spectrum
on a restricted commitment image. Adopting \(\mathbb F_{3^{n_u}}\) would
require its own field-element encoding: ordinary integer casts reach only
the prime subfield, and the current wire-modulus convention is unchanged.

\subsection{Scope of the instantiated application}

Paper 2's constructed systematic profile fixes the concrete cellular model,
full layout, identity integer coefficient map, raw payload carrier and
intrinsic modular valid locus. The results proved here compare its three
independently defined native metrics, construct a nonempty capacity-far
original-payload family and analyze its challenges. The calibration is
packaged through the existing
source-to-source and source-to-code interfaces; the existing near projection
uses the declared local scale \(2d_C\). The public-observation adapter retains
the original artifact and its actual extracted pair.

An application to a different maintained sheaf must prove the corresponding
source and coefficient identifications. An independently chosen HCP/frame
score needs its own upstream comparison; the native modular theorem does not
identify that score with edit distance. The global canonical-lift and
fixed-challenge obstructions constrain such stronger claims. The restricted
Euclidean theorem and the diagonal fixed-mix example retain their stated
profiles.

The selected profile gives a concrete relation for comparison with existing
commitment interfaces. In its complete identity layout, reduce the original
integer payload to \(v\in\mathbb F_q^{2n}\). Modular validity is the actual
coboundary-image relation, proved equivalent to paired RS membership. For a
fixed challenge, \(v\mapsto v_0+\gamma v_1\) is \(\mathbb F_q\)-linear.
This identification uses the chosen identity coefficient map; the earlier
generic interpretation \(a:\mathbb Z\to\mathbb Z\) is arbitrary.

Position-binding vector commitments \cite{catalanoFiore2011} and linear-map
vector commitments \cite{campanelliEtAl2022} concern authenticated coordinate
or function outputs. That authentication does not itself give the native
source-distance comparison above. A comparison with WHIR's constrained-RS
proximity interface \cite{arnonChiesaFenziYogev2024} additionally identifies
the polynomial constraints and fixed-pair challenge experiment. The fixed-mix
counterfamily prevents replacing this analysis by an arbitrary scalar fold.
Commitment/opening/verifier algorithms, metadata binding and their security
games belong to Paper 2. This paper proves the stated relation and metric
comparisons, without a new binding/hiding primitive, priority claim or
verification-cost advantage. These mathematical results do not require
completion of Paper 2's ledger, consensus protocol or Rust implementation.

\section{Formalization and Proof Ledger}

\subsection{Standalone hierarchy}

The machine-checked development separates generic coding mathematics from
application bridges:
\begin{leanledger}
\begin{description}[style=nextline]
  \item[\lean{Proximity/Core/BlockCode}]
    Relative distance, balls, and universal list-bound predicates.
  \item[\lean{Proximity/Core/ReedSolomon}]
    Evaluation-domain definitions, RS codes, and interleaving.
  \item[\lean{Proximity/Core/PlanePair}]
    Plane--pair--line equivalences.
  \item[\lean{Proximity/Core/FiniteThreshold}]
    Safe finite grids, greatest-safe and first-unsafe values, and adjacent
    safe/unsafe threshold certificates.
  \item[\lean{Proximity/Core/Distance} and \lean{Proximity/Core/RadiusShell}]
    Attained minimum distances, closed and strict nearness, and the exact
    half-open real intervals between consecutive Hamming-count shells.
  \item[\lean{Proximity/Capacity/Parameters}]
    Explicit \(\epsstar\) and field-size inequalities.
  \item[\lean{Proximity/Capacity/Realization} and \lean{Proximity/Capacity/ListPhase}]
    Actual Galois fields, exact-order subgroup domains for every family
    member, and instantiated list/MCA grid and real-radius conclusions.
  \item[\lean{Proximity/ListDecoding/Capacity}]
    Agreement-set injection, the matching received-word construction, the
    universal endpoint equality, the fractional-target equivalence, and the
    next-shell coefficient-vector injection.
  \item[\lean{Proximity/ListDecoding/PhaseSplit}]
    All three fractional-target regimes, the exact endpoint-safe real
    interval, its excluded supremum, and the largest safe grid count.
  \item[\lean{Proximity/MCA/Definitions}]
    Arbitrary-block-code same-set MCA definitions at count and relative
    radii.
  \item[\lean{Proximity/MCA/Threshold}]
    Exact worst bad-count spectrum and the equivalence between operational
    real error at \(r/n\) and the integer \(2^{-\lambda}\) budget.
  \item[\lean{Proximity/MCA/SpectrumCertificate}]
    Supplied-code grid-threshold pinning from a uniform upper bound at one
    radius and a concrete over-budget affine line at the next radius.
  \item[\lean{Proximity/MCA/SmallRadius}]
    Universal RS unit-spike and two-spike certificates, including
    \(B_{\mathcal C}(0)=1\) for every proper RS code.
  \item[\lean{Proximity/MCA/BoundaryField}]
    Primality, order-\(4096\) smooth domain, and the exact 128-bit radius
    transition for \(q=2^{128}+192513\).
  \item[\lean{Proximity/MCA/RealThreshold}]
    Count/real same-set event correspondence and the boundary-field safe
    interval \([0,1/4096)\), including supremum and nonattainment.
  \item[\lean{Proximity/MCA/Pair}]
    RS compressed interpolation certificates,
    bad-challenge injection, the uniform cardinality bound, and the
    explicit-family error corollary.
  \item[\lean{Proximity/Application/Calibration}]
    Generic source-validity model, pointwise distance non-collapse, and the
    uniform source-far to target-far implication.
  \item[\lean{Proximity/Application/DistanceCalibration}]
    Identification with attained scalar/pair code distance, interpolation
    covering bounds, and good-challenge farness with its exact boundary.
  \item[\lean{Proximity/Core/CanonicalLift} and \lean{Proximity/Core/Systematic}]
    Canonical real-lift spanning, actual systematic interpolation and
    integer/real syndrome compatibility. These are proved algebraic facts.
  \item[\lean{Proximity/Core/InterleavedMetric}]
    Native scalar-coordinate and whole-pair Hamming comparisons, retaining
    their different normalizations.
  \item[\lean{Proximity/Core/MetricSharpness}]
    The degree-\(k\) first column with zero second column attains factor two
    between pair and scalar-coordinate distance for proper RS codes.
  \item[\lean{Proximity/Core/DeepHole} and \lean{Proximity/Core/FixedChallenge}]
    Exact scalar/pair/flattened capacity distances for monomial received words,
    and capacity-far counterexamples to arbitrary fixed-challenge compression.
  \item[\lean{Proximity/MCA/ZeroDirection} and \lean{Proximity/MCA/Diagonal}]
    Explicit all-good and single-cancelling-challenge profiles, with their
    exact same-set bad events.
\end{description}
\end{leanledger}
The standalone import closure does not depend on HCP, protocol, consensus, or
runtime axiom bundles.
Both directions of the capacity-grid equality and its exact fractional-target
criterion are theorems, not axioms.  The next-shell lower bound is also
machine-checked.  The pair-MCA theorem and its \(2^{-128}\) family
instantiation are machine-checked as well.
The generic family corollaries retain their parameterized forms. The
realization module additionally constructs \(\mathbb F_{3^{n_u}}\), an
exact-order \(n_u\) root, and its injectively indexed subgroup for every
\(u\), and instantiates those corollaries without a realization hypothesis.
All source modules in this ledger live under \lean{lean/Proximity};
\lean{lean/ProximityCoding.lean} is only the standalone aggregate import target,
not a second source hierarchy.
The detailed object dictionary, theorem correspondence, construction ownership
and remaining mathematical extensions are recorded in
\lean{papers/paper3/SEMANTIC_PARITY.md}.

\subsection{Auxiliary source and projective proofs}

The optional geometric source in
\cref{sec:projective-observation-boundary} uses the existing
\lean{Proximity/Concrete/GrassmannianPlucker/Reconstruction} module.
Its \lean{recoverSubspace_frame} and \lean{recoverHarmonicProjector_frame}
theorems recover the full subspace and its fixed-metric projector.
\lean{TCB/PluckerObservation.recover_eq_operativeHarmonicComponent}
adds the explicit operative-kernel identification.
\lean{Proximity/Concrete/ProjectionChain/Core} composes the supplied
projective, cochain, plane and RS maps and their validity/near comparisons.
The exact-component obstruction in
\cref{prop:exact-component-obstruction} uses the existing
\lean{Sheaf/ObservationQuotient} proofs
\lean{HodgeEncryptedCochain.not_residual_norm_factors_exactComponent} and
\lean{HodgeEncryptedCochain.not_far_preserving_exactComponent}.
They retain the unrestricted Hodge carrier, a nonzero coexact direction,
exact-component factorization and zero target distance at the zero image.
The canonical systematic sheaf has \(\delta_1=0\), so this nonzero-coexact
premise is unavailable there.

These source-obstruction and projective results are canonical auxiliary proofs
outside the configured 52-target coding and 22-target native-modular
application selections. They do not enter the standalone coding import
closure; moving a statement into this manuscript does not add a selected
target. Their formal owners and the independent source metrics remain
distinct from Paper 2's codec, visibility and authentication contracts.
The projection-chain composition gives validity and near-side handoffs;
source-far preservation requires the separately stated calibration.

\subsection{Status ledger}

\begingroup
\small
\begin{longtable}{@{}>{\raggedright\arraybackslash}p{0.42\linewidth}
  >{\raggedright\arraybackslash}p{0.20\linewidth}
  >{\raggedright\arraybackslash}p{0.28\linewidth}@{}}
\toprule
Obligation & Status & Required endpoint\\
\midrule
\endfirsthead
\toprule
Obligation & Status & Required endpoint\\
\midrule
\endhead
Count/relative ball correspondence & Lean-proved & Equality at \(r/n\),
  for \(n>0\)\\
Plane--pair equivalence & Lean-proved & Bijection and Hamming isometry\\
Plane--affine-line equivalence & Lean-proved & Formal inverse maps\\
Smooth large-field family & Lean-proved & Actual fields and subgroup domains\\
Capacity-grid list equality & Lean-proved & Universal \(L_m=\binom nk\)\\
Endpoint fractional criterion & Lean-proved & Target iff
  \(\binom nk\le\epsstar|\F|\)\\
Next-shell list lower bound & Lean-proved & Explicit \(|\F|^m\) injection\\
Complete list phase split & Lean-proved & Three regimes and exact real interval\\
Generic minimum-distance interpretation & Lean-proved & Attained minima,
  covering and good-challenge farness\\
Pair-MCA certificate bound & Lean-proved &
  \(|\Bad_{\MCA}|\le\binom n{k+1}\)\\
MCA target-fraction inequality & Lean-proved &
  \(\binom n{k+1}\le\epsstar|\F|\) for the explicit family\\
Exact supplied-code MCA grid API & Lean-proved &
  Operational target iff \(2^{128}B_{\mathcal C}(r)\le|\F|\), with
  concrete next-line pinning\\
Exact RS small-radius certificates & Lean-proved &
  \(B_{\mathcal C}(0)=1\), and \(B_{\mathcal C}(1)\ge2\) for
  \(1\le k,\ k+2\le n\)\\
Smooth boundary-field instance & Lean-proved &
  Prime \(q=2^{128}+192513\equiv1\pmod {4096}\)\\
Boundary-field MCA threshold & Lean-proved &
  Greatest safe count \(0\); first unsafe count \(1\)\\
Explicit-family MCA greatest radius & Lean-proved &
  Full count radius \(n\) for the constructed family\\
Boundary-field real supremum & Lean-proved &
  \(1/4096\), not attained\\
Native scalar/pair metric comparison & Lean-proved &
  \(d_C\le d_P\le2d_C\)\\
Intrinsic modular application & Lean-proved & Actual integral coboundary
  validity, full original payload and \(d_C\le d_P=D_2\le2d_C\)\\
Sharp near scaling & Lean-proved & Factor two attained by \((u,0)\)\\
Canonical raw payload & Lean-proved & All three distances \(3/4\) at
  \((q,n,k)=(12289,4096,1024)\)\\
Monomial capacity far family & Lean-proved &
  Scalar, pair and flattened distances \((n-k)/n\)\\
Explicit challenge profiles & Lean-proved &
  All-good zero direction; diagonal bad set \(\{-1\}\)\\
Arbitrary fixed-mix obstruction & Lean-proved &
  Capacity-far pair with challenged word zero\\
Exact-component residual obstruction & Lean-proved, auxiliary &
  Nonfactorization and no positive far transfer under the stated
  unrestricted/nonzero-coexact hypotheses\\
Full projective reconstruction & Lean-proved, auxiliary &
  Subspace/projector recovery from genuine full-frame data\\
\bottomrule
\end{longtable}
\endgroup

The counting, calibration, boundary-field and canonical application proofs
use only the standard Lean logical axioms
\lean{propext}, \lean{Classical.choice}, and \lean{Quot.sound}.
The arithmetic certificates are kernel-checked: the boundary primality proof
uses an explicit Lucas factor chain, and modular powers are checked by binary
modular-exponentiation proof terms. The field, roots and domain parameters
are unchanged. Earlier versions used \lean{native_decide} for some of these
certificates; their dated dependency reports are historical evidence rather
than the trust footprint of the present proofs. This arithmetic change does
not construct a commitment protocol or a stronger source-metric comparison.

\subsection{Submission boundary}

A completed claim requires:
\begin{enumerate}[label=(\roman*)]
  \item no placeholder assumptions for the main list/MCA result;
  \item an explicit statement of imported literature theorems;
  \item all field-size and rounding conditions checked;
  \item a clear separation between proved theorems, computational evidence, and
    conjectures;
  \item a reproducible proof artifact and manuscript cross-reference.
\end{enumerate}

\subsection{Remaining interfaces}

The nine named standalone coding statements now have complete Lean
endpoints, including the all-parameter smooth-family realization, complete
list phase split, real safe intervals and their excluded suprema, and the
generic minimum-distance interpretation. The generic source-to-pair calibration implication
\[
  d_{\mathrm{source}}(x,\mathrm{valid})\ge\delta
  \Longrightarrow
  \Delta(\pi(x),\mathcal C^{\equiv2})\ge\delta
\]
is Lean-proved for an abstract distance model from its named pointwise
non-collapse field. The distance-calibration module identifies the target
model with the generic attained pair-code minimum. Pullback distance
makes that field reflexive; an application choosing a different pre-existing
source metric supplies one explicit comparison.
Separately, the intermediate list-decoding regime still requires new ideas if
one wants an optimal sub-endpoint radius for a supplied field that fails the
binomial target inequality.

The existing commitment-to-plane extraction, canonical plane/pair Hamming
isometry, exact transported-code distance equality, and public-observation
pullback calibration live in the concrete application hierarchy. They are
described in Section~7 and are outside the standalone
\lean{ProximityCoding} import closure. Section~7 supplies an explicit modular
instance of those interfaces. Comparison with an independently chosen HCP or
frame score remains a separate mathematical question.

The explicit source realization lives in the \lean{Sheaf/SystematicCode}
modules, including \lean{Modular}. Their actual integral cellular coboundary
is proved compatible with RS reduction. The storage modules
\lean{CoordinateCodec}, \lean{NoWrapCalibration} and \lean{ThreeLevelMetrics}
retain the original payload and separate the restricted Euclidean comparison
from the global native-metric comparison. The canonical specialization and
public-observation joins are in \lean{CanonicalSystematic} and
\lean{CanonicalThreeLevel}; \lean{CanonicalCalibration} composes the existing
interfaces with their declared near/far scales. These modules live under
\lean{Protocol/Storage}. Their file location does not change manuscript
ownership: Paper 2 specifies the shared construction and source/observation
contracts \cite{vanDerVlugt2026protocol}; this paper proves the coding and
quantitative metric comparisons for that exact construction. The two Paper 3
verification configurations separate core and bridge statements for source-size
purposes, without introducing another paper or protocol result.

These explicit profiles do not identify arbitrary old HCP presentations or
their frame scores with native modular edit distance. A stronger upstream
geometric comparison, or a different challenge/observation profile, needs its
own theorem. Such obligations remain separate from the completed generic
coding results and do not require Paper~2 consensus or Rust.


\bibliographystyle{alpha}
\bibliography{references}

\end{document}
